% Copyright 2020 by Robert Hildebrand
%This work is licensed under a
%Creative Commons Attribution-ShareAlike 4.0 International License (CC BY-SA 4.0)
%See http://creativecommons.org/licenses/by-sa/4.0/% 

\documentclass[../open-optimization/open-optimization.tex]{subfiles}
% 
\begin{document}

\chapter{Integer Programming Formulations}

\label{sec:IP-formulations} 
\begin{outcome}
\begin{enumerate}
\item[A.] Learn classic integer programming formulations.
\item[B.] Demonstrate different uses of binary and integer variables.
\item[C.]  Demonstrate the format for modeling an optimization problem with sets, parameters, variables, and the model.
\end{enumerate}
\end{outcome}
\begin{comment}
%We will learn about complexity theory in a subsequent section.  But take note of the complexity classification of each problem ( \polynomial,\npcomplete, \nphard) to get a sense of how difficult the problem is.  We will define these terms later.
\end{comment}

In this section, we will describe classical integer programming formulations.     These formulations may reflect  a real world problem exactly, or may be part of the setup of a real world problem.   

\section{Knapsack Problem}
The \emph{knapsack problem} can take different forms depending on if the variables are binary or integer.  The binary version means that there is only one item of each item type that can be taken.  This is typically illustrated as a backpack (knapsack) and some items to put into it (see  \autoref{fig:wiki/File/knapsack}), but has applications in many contexts.  

 \includefigurestatic[Knapsack Problem: which items should we choose take in the knapsack that maximizes the value while respecting the 15kg weight limit?][width=.4\linewidth][h]{wiki/File/knapsack}

\begin{general}{Binary Knapsack Problem}{\npcomplete}
Given an non-negative weight vector $\a \in \Q^n_+$, a capacity $b \in \Q_+$, and objective coefficients $\c \in \Q^n$, 
\begin{equation}
\begin{split}
\max \ \ & \c^\top \x\\
\st\ \ & \a^\top \x \leq b\\
& \x \in \{0,1\}^n
\end{split}
\end{equation}
\end{general}
 

 
\begin{figure}[H]
\begin{center}
\includegraphics[scale = 0.2]{knapsack}\footnotemark 
\end{center}
\label{fig:knapsack}
\caption{Knapsack Problem: which items should we choose take in the knapsack that maximizes the value while respecting the 15kg weight limit?}
\end{figure}
\footnotetext{\url{https://en.wikipedia.org/wiki/Knapsack_problem}}

\begin{examplewithcode}{Knapsack}{\href{https://github.com/open-optimization/open-optimization-or-examples/blob/master/integer-programming/knapsack-problem.ipynb}{Gurobipy}}
\label{example:knapsack}
You have a knapsack (bag) that can only hold W = 15 kgs.  There are 5 items that you could possibly put into your knapsack.  The items (weight, value) are given as:
(12 kg, $\$$4), (2 kg, $\$$2), (1kg, $\$$2), (1kg, $\$$1), (4kg, $\$$10).  Which items should you take to maximize your value in the knapsack? See \autoref{fig:wiki/File/knapsack}.\\

\modelhead{Variables}
\begin{itemize}
\item let $x_i = 0$ if item $i$ is in the bag
\item let $x_i = 1$ if item $i$ is not in the bag
\end{itemize}
\modelhead{Model}
\begin{align}
\max \quad & 4 x_1 + 2 x_2 + 2 x_3 + 1 x_4 + 10 x_5   \tag{Total value} \\
\st        & 12 x_1 + 2 x_2 + 1 x_3 + 1 x_4 + 4 x_5 \leq 15   \tag{Capacity bound} \\
& x_i \in \{0,1\} &&\quad \text{for all } i=1, \dots, 5   \tag{Item taken or not}
\end{align}
\end{examplewithcode}
In the integer case, we typically require the variables to be non-negative integers, hence we use the notation $\x \in \Z^n_+$.  This setting reflects the fact that instead of single individual items, you have item types of which you can take as many of each type as you like that meets the constraint.
\begin{general}{Integer Knapsack Problem}{\npcomplete}
Given an non-negative weight vector $\a \in \Q^n_+$, a capacity $b \in \Q_+$, and objective coefficients $c \in \Q^n$, 
\begin{equation}
\begin{split}
\max \ \ & \c^\top \x\\
\st\ \ & \a^\top \x \leq b\\
& \x \in \Z^n_+
\end{split}
\end{equation}
\end{general}
We can also consider an equality constrained version
\begin{general}{Equality Constrained Integer Knapsack Problem}{\nphard}
Given an non-negative weight vector $\a \in \Q^n_+$, a capacity $b \in \Q_+$, and objective coefficients $c \in \Q^n$, 
\begin{align}
\max \quad & \c^\top \x \\
\st        & \a^\top \x = b \\
& \x \in \Z^n_+
\end{align}
\end{general}

\begin{examplewithcode}{Min Coins}{\href{https://github.com/open-optimization/open-optimization-or-examples/blob/master/integer-programming/coins-minimize-problem.ipynb}{Gurobipy}}
\label{ex:min-coins}
Using pennies, nickels, dimes, and quarters, how can you minimize the number of coins you need to to make up a sum of $83\cent$? 

\modelhead{Variables}
\begin{itemize}
\item Let $p$ be the number of pennies used
\item Let $n$ be the number of nickels used
\item Let $d$ be the number of dimes used
\item Let $q$ be the number of quarters used
\end{itemize}
\modelhead{Model}
\begin{align*}
\min \quad & p + n + d + q   \tag{total number of coins used} \\
\st        & p + 5n + 10d + 25 q = 83 & \text{sums to } 83 \cent \\
& p,d,n,q \in \Z_+   \tag{each is a non-negative integer}
\end{align*}
\end{examplewithcode}

\section{Capital Budgeting}

\input{optimization/open-optimization-examples/example-capital-budgeting.tex}

\section{Capacitated Lot Sizing Problem (CLSP)}

The \textit{Capacitated Lot Sizing Problem (CLSP)} is a fundamental optimization problem in production planning, where the goal is to determine an optimal production schedule that minimizes total costs while adhering to capacity constraints. The objective function accounts for production costs, setup costs, and inventory holding costs over a given planning horizon.

Let \( T \) denote the number of discrete time periods, indexed by \( t \in \{1,2,\dots,T\} \). The problem parameters are defined as follows:
\begin{itemize}
    \item \( d_t \) is the demand in period \( t \),
    \item \( p_t \) represents the unit production cost,
    \item \( f_t \) is the fixed setup cost incurred if production takes place,
    \item \( h_t \) denotes the unit inventory holding cost, and
    \item \( c_t \) represents the maximum production capacity in period \( t \).
\end{itemize}

The decision variables are:
\begin{itemize}
    \item \( x_t \) denotes the quantity produced in period \( t \),
    \item \( s_t \) is the inventory level at the end of period \( t \), and
    \item \( y_t \) is a binary variable that takes value 1 if production occurs in period \( t \), and 0 otherwise.
\end{itemize}

The CLSP can be formulated as the following mixed-integer linear program:
\begin{align}
\min \quad & \sum_{t=1}^{T} \left( p_t x_t + f_t y_t + h_t s_t \right)   \tag{Minimize total cost}\label{eq:objective} \\
\st        & s_{t-1} + x_t - d_t = s_t &&\quad \text{for all } t = 1,2,\dots,T   \tag{Inventory balance}\label{eq:inventory} \\
    & x_t \leq y_t c_t &&\quad \text{for all } t = 1,2,\dots,T   \tag{Capacity and setup constraint}\label{eq:capacity} \\
    & x_t, s_t \geq 0 &&\quad \text{for all } t = 1,2,\dots,T   \tag{Non-negativity constraints}\label{eq:nonnegativity} \\
    & y_t \in \{0,1\} &&\quad \text{for all } t = 1,2,\dots,T   \tag{Binary production decision}\label{eq:binary}
\end{align}

The objective function \eqref{eq:objective} minimizes the total cost over the planning horizon, incorporating production, setup, and inventory holding costs. The inventory balance constraint \eqref{eq:inventory} ensures that demand in each period is met using either inventory from the previous period or newly produced units. The constraint \eqref{eq:capacity} enforces production capacity limits and ensures that a setup cost is incurred when production occurs. Constraints \eqref{eq:nonnegativity} enforce non-negativity on production and inventory levels, while \eqref{eq:binary} ensures that the setup decision is binary.

For simplicity, we assume an initial inventory level of zero.

\section{Facility Location}

The basic model of the facility location problem is to determine where to place your stores or facilities in order to be close to all of your customers and hence reduce the costs transportation to your customers.  Each customer is known to have a certain demand for a product, and each facility has a capacity on how much of that demand it can satisfy.  Furthermore, we need to consider the cost of building the facility in a given location.

This basic framework can be applied in many types of problems and there are a number of variants to this problem.   We will address two variants: the \emph{capacitated facility location problem} and the \emph{uncapacitated facility location problem}.  

\subsection{Capacitated Facility Location}
We present 3 ways to model this problem.

\begin{general}{Capacitated Facility Location 1}{\npcomplete}
Given connection costs \( c_{ij} \), fixed building costs \( f_i \), demands \( d_j \), and facility capacities \( u_i \), the \textbf{capacitated facility location} problem is formulated as follows:

\modelhead{Sets}
\begin{itemize}
\item Let \( I = \{1,\dots, n\} \) be the set of facilities.
\item Let \( J = \{1, \dots, m\} \) be the set of customers.
\end{itemize}

\modelhead{Parameters}
\begin{itemize}
\item \( f_i \) — the cost of opening facility \( i \).
\item \( c_{ij} \) — the cost of serving one unit of demand of customer \( j \) from facility \( i \).
\item \( u_i \) — the capacity (in units) of facility \( i \).
\item \( d_j \) — the total demand (in units) of customer \( j \).
\end{itemize}

\modelhead{Variables}
\begin{itemize}
\item \( y_i \in \{0,1\} \) — equals 1 if facility \( i \) is opened; 0 otherwise.
\item \( x_{ij} \geq 0 \) — the number of units of customer \( j \)'s demand served by facility \( i \).
\end{itemize}

\modelhead{Model}
\begin{align}
\min \quad & \sum_{i=1}^n \sum_{j=1}^m c_{ij} x_{ij} + \sum_{i=1}^n f_i y_i   \tag{Total cost} \\
\st        & \sum_{i=1}^n x_{ij} = d_j &&\quad \text{for all } j = 1,\dots,m   \tag{Meet demand} \\
& \sum_{j=1}^m x_{ij} \leq u_i y_i &&\quad \text{for all } i = 1,\dots,n   \tag{Facility capacity} \\
& x_{ij} \geq 0 &&\quad \text{for all } i = 1,\dots,n,\; j = 1,\dots,m   \tag{Nonnegativity} \\
& y_i \in \{0,1\} &&\quad \text{for all } i = 1,\dots,n   \tag{Open facility}
\end{align}

\end{general}

Next we present a method that models it tracking the fraction of demand served.
\begin{general}{Capacitated Facility Location 2}{\npcomplete}
Given costs connections $c_{ij}$ and fixed building costs $f_i$, demands $d_j$ and capacities $u_i$, the capacitated facility location problem is 

\modelhead{Sets}
\begin{itemize}
\item Let $I = \{1,\dots, n\}$ be the set of facilities.
\item Let $J = \{1, \dots, m\}$ be the set of customers.
\end{itemize}

\modelhead{Parameters}
\begin{itemize}
\item $f_i$ - the cost of opening facility $i$.
\item $c_{ij}$ - the cost of fulfilling the complete demand of customer $j$ from facility $i$.
\item $u_{i}$ - the capacity of facility $i$.
\item $d_{j}$ - the demand by customer $j$.
\item $c_j$ is given in column "$\mathbb{E}$[Revenue]".
\item $b_i$ is given in row "Resources available".
\item $a_{ij}$ given in row $j$, and column for week $i$.
\end{itemize}

\modelhead{Variables}
\begin{itemize}
\item Let 
\begin{equation*}
y_{i} = \begin{cases}
1 & \text{if we open facility $i$,}\\
0 & \text{otherwise.}
\end{cases}
\end{equation*}
\item Let $x_{ij} \geq 0$ be the fraction of demand of customer $j$ satisfied by facility $i$.
\item let $x_i = 1$ if investment $i$ is not chosen
\end{itemize}

\modelhead{Model}
\begin{align}
\min \quad & \displaystyle\sum_{i=1}^n\sum_{j=1}^mc_{ij}x_{ij}+\sum_{i=1}^nf_iy_i   \tag{total cost} \\
\st        & \displaystyle\sum_{i=1}^nx_{ij}=1 &&\quad \text{for all } j=1,\dots,m   \tag{assign demand to facility} \\
& \displaystyle \sum_{j=1}^md_jx_{ij}\leqslant u_iy_i &&\quad \text{for all } i=1\dots,n   \tag{capacity of facility $i$} \\
& x_{ij}\geqslant0 &&\quad \text{for all } i=1,\dots,n \text{ and }j=1,\dots,m   \tag{nonnegative fraction of demand satisfied} \\
& y_i\in\{0,1\} &&\quad \text{for all } i=1,\dots,n   \tag{open/not open facility}
\end{align}

\end{general}

Finally, we present a version that has the same variables, but the constraints model the feasible solutions in a slightly different way.

\begin{align*}
\min \quad & \sum_{i = 1}^m\sum_{j=1}^n c_{ij}x_{ij}+\sum_{i = 1}^mf_iy_i \\
\st        & \sum_{i = 1}^mx_{ij} = 1,~~~~~j =1,\ldots n \\
& \displaystyle \sum_{j=1}^md_jx_{ij}\leqslant u_i &&\quad \text{for all } i=1\dots,n   \tag{capacity of facility $i$} \\
& x_{ij}\leq y_i,~~~~~ i = 1,\ldots m,~~~ j = 1,\ldots,n \\
& x_{ij} \geq 0~~~~~ i = 1,\ldots,m, ~~~j = 1,\ldots,n \\
& y_i\in \{0,1\}, \  \  i = 1,\ldots,m
\end{align*}

\begin{examplewithcode}{Capacitated Facility Location: Retail Distribution Example}{\href{https://github.com/open-optimization/open-optimization-or-examples/blob/master/integer-programming/facility-location-capacitated.ipynb}{Gurobipy}}

\textbf{Context:} A retail company plans to establish distribution centers across a country to serve its stores efficiently. The company faces decisions on which distribution centers to open and which stores each center should serve. Costs associated with opening each center and serving stores from them are known, and each center has a maximum capacity. Additionally, each store has a known demand.

\textbf{Given Data:} (\href{https://github.com/open-optimization/open-optimization-or-examples/blob/master/integer-programming/distribution-data.xlsx}{distribution-data.xlsx})

\begin{itemize}
    \item Number of potential distribution centers ($n$): 3
    \item Number of stores ($m$): 4
\end{itemize}

\begin{table}[ht]
\centering
\caption{Combined Data for Distribution Centers and Stores}
\begin{tabular}{lcc|cccc}
\toprule
 & \(\mathbf{f_i}\) & \(\mathbf{u_i}\) & \multicolumn{4}{c}{\(\mathbf{c_{ij}}\) (Cost to Serve a Store)} \\
\cmidrule(lr){4-7}
\textbf{Center} & (\$) & (units) & Store 1 & Store 2 & Store 3 & Store 4 \\
\midrule
Center 1 & 150,000 & 100 & \$50 & \$60 & \$70 & \$85 \\
Center 2 & 100,000 & 80  & \$75 & \$45 & \$55 & \$50 \\
Center 3 & 180,000 & 90  & \$65 & \$80 & \$40 & \$60 \\
\midrule
\textbf{Store Demand \(d_j\)} & & & 400 & 300 & 200 & 250 \\
\bottomrule
\end{tabular}
\end{table}

\textbf{Model Formulation:}
The capacitated facility location model described earlier can be applied to this scenario with the given data to determine the optimal number and location of distribution centers to open, and which stores each center should serve.
\begin{center}
\includegraphics[scale = 0.75]{Figures/facility-location}
\end{center}
\end{examplewithcode}

\subsection{Uncapacitated Facility Location}
\begin{general}{Uncapacitated Facility Location}{\npcomplete}
Given costs connections $c_{ij}$ and fixed building costs $f_i$, the uncapacitated facility location problem is 
\begin{equation}
\begin{array}{rl}
\min & \displaystyle\sum_{i=1}^n\sum_{j=1}^mc_{ij}z_{ij}+\sum_{i=1}^nf_ix_i \\
\st & \displaystyle\sum_{i=1}^nz_{ij}=1 \text{ for all }j=1,\dots,m \\
& \displaystyle \sum_{j=1}^mz_{ij}\leqslant Mx_i\text{ for all }i=1\dots,n \\
&z_{ij}\in\{0,1\}\text{ for all }i=1,\dots,n \text{ and }j=1,\dots,m\\
&x_i\in\{0,1\}\text{ for all } i=1,\dots,n
\end{array}
\end{equation}
Here $M$ is a large number and can be chosen as $M = m$, but could be refined smaller if more context is known.
\end{general}

Alternative model!

\begin{equation}
\begin{array}{rl}
\min & \displaystyle\sum_{i=1}^n\sum_{j=1}^mc_{ij}z_{ij}+\sum_{i=1}^nf_ix_i \\
\st & \displaystyle\sum_{i=1}^nz_{ij}=1 \text{ for all }j=1,\dots,m \\
& z_{ij}\leqslant x_i\text{ for all }i=1\dots,n \text{ for all } j = 1, \dots, m \\
&z_{ij}\in\{0,1\}\text{ for all }i=1,\dots,n \text{ and }j=1,\dots,m\\
&x_i\in\{0,1\}\text{ for all } i=1,\dots,n
\end{array}
\end{equation}
\newpage

\section{Set Covering}

The \emph{set covering} problem can be used for a wide array of problems.    We will see several examples in this section.

\begin{general}{Set Covering}{\npcomplete}
\label{general:set-covering}
Given a set $V$ with subsets $V_1, \dots, V_l$, determine the smallest subset $S \subseteq V$ such that 
$S \cap V_i \neq \emptyset$ for all $i=1, \dots, l$.

The set cover problem can be modeled as
\begin{equation}
\begin{split}
\min \ \ & \one^\top x\\
\st \ \ & \sum_{v \in V_i} x_v \geq 1 \text{ for all } i =1, \dots, l \\ 
& x_v \in \{0,1\} \text{ for all } v \in V
\end{split}
\end{equation}
where $x_v$ is a 0/1 variable that takes the value $1$ if we include item $j$ in set $S$ and $0$ if we do not include it in the set $S$.  
\end{general}

%
\begin{examplewithcode}{Capital Budgeting}{https://github.com/open-optimization/open-optimization-or-examples/blob/master/integer-programming/capital-budgeting.ipynb}
\modelhead{Sets}
\begin{itemize}
\item Let $I = \{1,2\}$ be the set of time periods.
\item Let $J = \{1, 2,3\}$ be the set of possible investments.
\end{itemize}
%
\modelhead{Parameters}
\begin{itemize}
\item $c_j$ is given in column "$\mathbb{E}$[Revenue]".
\item $b_i$ is given in row "Resources available".
\item $a_{ij}$ given in row $j$, and column for week $i$.
\end{itemize}
%
\modelhead{Variables}
\begin{itemize}
\item let $x_i = 0$ if investment $i$ is chosen
\item let $x_i = 1$ if investment $i$ is not chosen
\end{itemize}
%
%The explict model is given by\\
\modelhead{Model}
\begin{align*}
\max \quad & 10x_1+8x_2+6x_3   \tag{Total Expected Revenue} \\
%	\st \ \ & 
%	3x_1+1x_2+2x_3\leq5 \tag{ Resource constraint week 1}\\
%	&4x_1+2x_2+1x_3\leq6 \tag{ Resource constraint week 2}\\
%&x_j \in \{0,1\},\  j = 1,2,3
\end{align*}
%
\end{examplewithcode}
%

One specific type of set cover problem is the \emph{vertex cover} problem.
\begin{general}{Example: Vertex Cover}{\npcomplete}
\label{ex:vertex-cover}
Given a graph $G = (V,E)$ of vertices and edges, we want to find a smallest size subset $S \subseteq V$ such that every for every $e = (v,u) \in E$, either $u$ or $v$ is in $S$.   

We can write this as a mathematical program in the form:
\begin{equation}
\begin{split}
\min \ \ & \one^\top x\\
\st \ \ & x_u + x_v \geq 1 \text{ for all } (u,v) \in E \\ 
& x_v \in \{0,1\} \text{ for all } v \in V.
\end{split}
\end{equation}
\end{general}

\begin{examplewithcode}{Set cover: Fire station placement}{https://github.com/open-optimization/open-optimization-or-examples/blob/master/integer-programming/fire-station-covering.ipynb}
\label{example:fire-station}

In the fire station problem, we seek to choose locations for fire stations such that any district either contains a fire station, or neighbors a district that constains a fire station.  Figure~\ref{fig:tikz/Illustration1.pdf} depicts the set of districts and an example placement of locations of fire stations.  How can we minimize the the total number of fire stations that we need?

\modelhead{Sets}
\begin{itemize}
\item Let $V$ be the set of districts ($V = \{1, \dots, 16\}$)
\item Let $V_i$ be the set of districts that neighbor district $i$ (e.g. $V_1 = \{2,4,5\}$).
\end{itemize}

\modelhead{Variables}
\begin{itemize}
\item let $x_i = 1$ if district $i$ is chosen to have a fire station.
\item let $x_i = 0$ otherwise.
\end{itemize}
\modelhead{Model}
\begin{align}
\min \quad & \sum_{i \in V} x_i \tag{ \# \text{ open fire stations}} \\
\st        & x_i + \sum_{j \in V_i} x_j \geq 1 &&\quad \text{for all } {i \in V}   \tag{Station proximity requirement} \\
& x_i \in \{0,1\} &&\quad \text{for all } i\in V   \tag{station either open or closed}
\end{align}
\end{examplewithcode}

\includefiguresource[Layout of districts and possible locations of fire stations.][][h]{tikz/Illustration1.pdf}
\includefiguresource[Set cover representation of fire station problem.  For example, choosing district 16 to have a fire station covers districts 13, 15, and 16.]{tikz/Illustration2.pdf}
\includefiguresource[Graph representation of fire station problem.  Every node is connected to a chosen node by an edge]{tikz/Illustration3.pdf}

\input{optimization/open-optimization-examples/example-set-covering.tex}

\subsection{Covering (Generalizing Set Cover)}
We could also allow for a more general type of set covering where we have non-negative integer variables and a right hand side that has values other than $1$.
\begin{general}{Covering}{\npcomplete}
\label{general:covering}
Given a non-negative matrix $A \in \Z_+^{m \times n}$, a non-negative vector $b \in \Z^m$, and an objective vector $c \in \R^n$, the set cover problem is
\begin{equation}
\begin{split}
\max \ \ & c^\top x\\
\st. \ \ & Ax \geq b \\
& x \in\Z^n_+.
\end{split}
\end{equation}
\end{general}

%
\begin{examplewithcode}{Nurse Scheduling}{code:nurse-scheduling}
%Suppose that we have $n$ full-time nurses employed and $p$ part-time nurses employed.  We need to schedule nurses such that, for each day of the week, we have enough nurses on hand to cover all the tasks.  
%
\end{examplewithcode}

\section{Graph Coloring}

Graph coloring \refincludefigurestatic{wiki/File/Petersen-graph-3-coloring.png} is 
is problem of finding a minimum coloring in a graph can be formulated in many ways. Each vertex is assigned a color and two vertices cannot share the same color if they are connected by an edge.

In order to present integer programming formulations, we use $x_{i j}$ to denote binary variables, with $i \in V$ and $1 \leqslant j \leqslant n$, where $x_{i j}=1$ if color $j$ is assigned to vertex $i$ and $x_{i j}=0$ otherwise. We also define $n$ binary variables $w_j$ for $j=1, \ldots, n$, that indicate whether color $j$ is used in some node, i.e., $w_j=1$ if $x_{i j}=1$ for some vertex $i$.
The following is a classical integer programming formulation:
$$
\begin{aligned}
\min & \sum_{j=1}^n w_j \\
& \sum_{j=1}^n x_{i j}=1 \quad \forall i \in V, \\
& x_{i j}+x_{k j} \leqslant w_j \quad \forall\{i, k\} \in E, \quad 1 \leqslant j \leqslant n, \\
& x_{i j} \in\{0,1\} \quad \forall i \in V, 1 \leqslant j \leqslant n\\
& w_j \in\{0,1\} \quad 1 \leqslant j \leqslant n .
\end{aligned}
$$

We can improve upon this model.  We suggest two possible models from~\cite{MENDEZDIAZ2008159}.  Please read their paper for more advanced models with additional families of cuts added.

\textbf{ Color order model:}. 
Given a $(k)$-coloring, selecting any $k$ colors from the set $\{1, \ldots, n\}$ results in a valid solution, and these solutions are deemed equivalent. To reduce redundancy among these solutions, we stipulate that color $j$ can only be allocated to a vertex if color $j-1$ has previously been assigned. By doing so, we exclude symmetrical $(k)$-colorings that use colors with labels exceeding $k$. To define these valid solutions, we simply need to incorporate the subsequent set of constraints:
$$
\begin{aligned}
& w_j \leqslant \sum_{i \in V} x_{i j} \quad \forall 1 \leqslant j \leqslant n, \\
& w_j \geqslant w_{j+1} \quad \forall 1 \leqslant j \leqslant n-1 .
\end{aligned}
$$
The reduction of feasible solutions is very significant.

\textbf{ Independent set order model:} 
In a $(k)$-coloring, permutations of the initial $k$ colors lead to equivalent outcomes. To reduce the occurrence of these similar solutions, the subsequent model introduces stricter constraints than the previous one. It mandates that the count of vertices colored with $j$ should be at least equal to or surpass the count of vertices colored with $j+1$. We incorporate the ensuing inequalities:
$$
\begin{aligned}
& w_j \leqslant \sum_{i \in V} x_{i j} \quad \forall 1 \leqslant j \leqslant n \\
& \sum_{i=1}^n x_{i j} \geqslant \sum_{i=1}^n x_{i j+1} \quad \forall 1 \leqslant j \leqslant n-1 .
\end{aligned}
$$

As an exercise, try implementing these different models and compare the solve times.

\newpage
\section{Basic Modeling Tricks - Using Binary Variables}

In this section, we describe ways to model a variety of constraints that commonly appear in practice.  The goal is changing constraints described in words to constraints defined by math.

Binary variables can allow you to model many types of constraints.  We discuss here varios logical constraints where we assume that $x_i \in \{0,1\}$ for $i=1, \dots, n$.  We will take the meaning of the variable to be selecting an item.

\begin{enumerate}
 \item If item $i$ is selected, then item $j$ is also selected.
 \begin{equation}
 x_i \leq x_j
 \end{equation}
 \begin{enumerate}
 \item If any of items $1, \dots, 5$ are selected, then item $6$ is selected.
 \begin{equation}
 x_1 + x_2 + \dots + x_5 \leq 5 \cdot x_6
 \end{equation}
 Alternatively!
 \begin{equation}
 x_i \leq x_6 \ \ \ \ \text{ for all } i=1, \dots, 5
 \end{equation}
 \end{enumerate}
 \item If item $j$ is not selected, then item $i$ is not selected.
 \begin{equation}
x_i \leq x_j
 \end{equation}
  \begin{enumerate}
 \item If item $j$ is not selected, then all items $1, \dots, i$ are not selected.
 \begin{equation}
 x_1 + x_2 + \dots + x_i \leq i \cdot x_j
 \end{equation}
 \end{enumerate}
  \item If item $j$ is not selected, then item $i$ is not selected.
 \begin{equation}
x_i \leq x_j
 \end{equation}
\item Either item $i$ is selected or item $j$ is selected, but not both.
 \begin{equation}
 x_i + x_j = 1
 \end{equation}
\item Item $i$ is selected or item $j$ is selected or both.
 \begin{equation}
 x_i + x_j \geq 1
 \end{equation}
\item If item $i$ is selected, then item $j$ is not selected.
 \begin{equation}
x_j \leq (1-x_i)
 \end{equation}

\item At most one of items $i ,j$, and $k$ are selected.
 \begin{equation}
 x_i + x_j + x_k \leq 1
 \end{equation}
\item At most two of items $i,j,$ and $k$ are selected.
 \begin{equation}
 x_i + x_j + x_k \leq 2
 \end{equation}
\item Exactly one of items $i,j,$ and $k$ are selected.
 \begin{equation}
  x_i + x_j + x_k = 1
 \end{equation}
\end{enumerate}

These tricks can be connected to create different function values.  

\begin{example}{Variable takes one of three values}{}
Suppose that the variable $x$ should take one of the three values $\{4, 8, 13\}$. This can be modeled using three binary variables as
\begin{align*}
x=4z_1 +8z_2 +13z_3\\
 z_1 +z_2 +z_3 =1\\
z_i \in \{0, 1\} \text{ for } i = 1, 2,3.
\end{align*}
As a convenient addition, if we want to add the possibility that it takes the value $0$, then we can model this as 
\begin{align*}
x=4z_1 +8z_2 +13z_3\\
 z_1 +z_2 +z_3 \leq 1\\
z_i \in \{0, 1\} \text{ for } i = 1, 2,3.
\end{align*}
\end{example}

We can also model variable increases at different amounts.  \begin{example}{Discount for buying more}{}
Suppose you can choose to buy 1, 2, or 3 units of a product, each with a decreasing cost.  The first unit is \$10, the second is \$5, and the third unit is \$3.
\begin{align*}
x=10z_1 +5z_2 +3z_3\\
 z_1\geq z_2 \geq z_3\\
z_i \in \{0, 1\} \text{ for } i = 1, 2,3.
\end{align*}
Here, $z_i$ represents if we buy the $i$th unit.  The inequality constraints impose that if we buy unit $j$, then we must by all units $i$ with $i < j$.
\end{example}

In this section, we describe ways to model a variety of constraints that commonly appear in practice.  The goal is changing constraints described in words to constraints defined by math.
\subsection{Big M constraints - Activating/Deactivating Inequalities}
Big M comes again!   It's extremely useful when trying to activate constrants based on a binary variable.

For intance, if we don't rent a bus, then we can have at most 3 passengers join us on our trip.   Consider passengers $A,B,C,D,E$ and let $x_i \in \{0,1\}$ be 1 if we take passenger $i$ and 0 otherwise.
We can model the constraint that we can have at most 5 passengers as 
$$
x_A + x_B + x_C + x_D + x_E \leq 3.
$$

We want to be able to activate this constraint in the event that we don't rent a bus.

Let $\delta \in \{0,1\}$ be 1 if rent a bus, and 0 otherwise.  

Then we want to say
\begin{quote}
If $\delta = 0$ , then 
$$
x_A + x_B + x_C + x_D + x_E \leq 3.
$$
\end{quote}

We can formulate this using a big-M constraint as

\begin{equation}
x_A + x_B + x_C + x_D + x_E \leq 3 + M \delta.
\end{equation}

Notice the two case
$$
\begin{cases}
x_A + x_B + x_C + x_D + x_E \leq 3  & \text{ if } \delta = 0\\
x_A + x_B + x_C + x_D + x_E \leq 3 +M & \text{ if } \delta = 1
\end{cases}
$$
In the second case, we choose $M$ to be so large, that the second case inequality is vacuous.    That said, choosing smaller $M$ values (that are still valid) will help the computer program solve the problem faster.  In this case, it suffices to let $M = 2$.

We can speak about this technique more generally as
\begin{general}{Big-M: If then}{}{}
We aim to model the relationship
\begin{equation}
\text{ If } \ \ \   \delta = 0, \ \ \  \text{then  } \ \ \ a^\top x \leq b.
\end{equation}
By letting $M$ be an upper bound on the quantity $a^\top x - b$, we can model this condition as 
\begin{equation}
\begin{split}
a^\top x - b& \leq M\delta\\
\delta & \in \{0,1\}
\end{split}
\end{equation}
\end{general}

Below are several applications of this modeling technique:

\begin{itemize}
  \item \textbf{Facility Location with Fixed Costs:}\\
  To ensure that a facility \(i\) must be open in order to serve customer \(j\), we write
  \[
  x_{ij} \le M \delta_i,
  \]
  where \(x_{ij}\) is the quantity shipped and \(\delta_i = 1\) if facility \(i\) is open.

\item \textbf{Piecewise Linear or Step Costs:}\\
%  Suppose a cost function is composed of linear segments. We can activate the constraints of the \(k\)-th segment only if binary \(\delta_k = 1\) via:
\[
%  y - (m_k x + b_k) \le M(1 - \delta_k).
\]

  \item \textbf{Power Plant Operation Constraints:}\\
  A generator can only produce power when it is turned on. Letting \(\delta_t = 1\) indicate that the generator is on at time \(t\), we write:
  \[
  P_{\min} \delta_t \le P_t \le P_{\max} \delta_t.
  \]

  \item \textbf{Task Scheduling with Precedence Constraints:}\\
  To ensure task \(B\) starts only after task \(A\) ends, if a binary \(\delta = 1\) indicates the precedence relation is active, we write:
  \[
  \text{Start}_B \ge \text{End}_A - M(1 - \delta).
  \]

  \item \textbf{Routing and Time Windows:}\\
  In vehicle routing problems, to model time consistency across arcs, we use:
  \[
  t_j \ge t_i + s_i - M(1 - \delta_{ij}),
  \]
  where \(\delta_{ij} = 1\) indicates travel from node \(i\) to node \(j\), and \(s_i\) is service time at node \(i\).

  \item \textbf{Production Planning and Setup Costs:}\\
  To ensure production \(x_t\) is positive only if a setup decision \(\delta_t = 1\) is made:
  \[
  x_t \le M \delta_t.
  \]
\end{itemize}
\newpage
\begin{example}{Network Flow with Toll Charges Using Big-M Constraints}{}
Consider a transportation network where some arcs incur a toll cost if used. Let:
\begin{itemize}
  \item \( G = (V, A) \) be a directed graph with nodes \(V\) and arcs \(A\).
  \item \( x_{ij} \geq 0 \) be the flow on arc \((i,j) \in A\).
  \item \( \delta_{ij} \in \{0,1\} \) indicate whether arc \((i,j)\) is used.
  \item \( c_{ij} \) be the per-unit flow cost on arc \((i,j)\).
  \item \( t_{ij} \) be the fixed toll incurred if arc \((i,j)\) is used.
  \item \( M_{ij} \) be a large constant upper bound on possible flow through arc \((i,j)\).
\end{itemize}

The objective is to route flow from a source node \(s\) to a sink node \(t\) at minimum total cost, including tolls when applicable:
\[
\min \sum_{(i,j) \in A} \left( c_{ij} x_{ij} + t_{ij} \delta_{ij} \right)
\]

\noindent Subject to:
\begin{align*}
  &\sum_{j: (i,j) \in A} x_{ij} - \sum_{j: (j,i) \in A} x_{ji} =
  \begin{cases}
    1 & \text{if } i = s \\
    -1 & \text{if } i = t \\
    0 & \text{otherwise}
  \end{cases} &&\text{(Flow conservation)} \\
  &x_{ij} \le M_{ij} \delta_{ij} &&\text{(Big-M: toll applies only if arc is used)} \\
  &x_{ij} \ge 0, \quad \delta_{ij} \in \{0,1\} &&\text{for all } (i,j) \in A
\end{align*}

\end{example}

\subsection*{NYC Transportation Network with Tolls}

Figure~\ref{fig:nyc_network} illustrates a simplified transportation network over the five boroughs of New York City, with selected connections to model bridges and tunnels. Each node represents a borough, and directed arcs represent feasible travel routes. Some of these arcs include tolls, which are incurred only in one direction, mimicking real-world tolling policies.

\paragraph{Nodes.}
The nodes in the graph correspond to the following boroughs:
\begin{itemize}
    \item Manhattan
    \item Brooklyn
    \item Queens
    \item Bronx
    \item Staten Island
\end{itemize}

\paragraph{Arcs.}
The directed arcs model travel between boroughs, with an optional toll charge. These arcs represent crossings such as bridges or tunnels, and are defined as follows:

\begin{center}
\begin{tabular}{lcc}
\toprule
\textbf{From} & \textbf{To} & \textbf{Toll} \\
\midrule
Manhattan & Brooklyn       & Yes (e.g., Brooklyn-Battery Tunnel) \\
Manhattan & Queens         & Yes (e.g., Queens-Midtown Tunnel) \\
Brooklyn  & Queens         & No  \\
Manhattan & Bronx          & Yes (e.g., RFK/Triborough Bridge) \\
Bronx     & Queens         & Yes \\
Manhattan & Staten Island  & Yes (e.g., Verrazzano-Narrows Bridge) \\
Staten Island & Brooklyn   & No  \\
\bottomrule
\end{tabular}
\end{center}

\paragraph{Toll Modeling.}
Each arc \((i, j)\) in the network has an associated binary toll flag. If a toll applies, the arc is subject to a fixed toll cost only if it is used. This can be modeled in a mixed-integer program using a binary variable \(\delta_{ij} \in \{0,1\}\) that indicates whether the arc is selected. A Big-\(M\) constraint of the form:
\[
x_{ij} \le M \cdot \delta_{ij}
\]
can be used to enforce toll cost and flow constraints, where \(x_{ij}\) is the flow over arc \((i,j)\), and \(M\) is a suitably large upper bound on the possible flow.

\paragraph{Legend.}
In the figure, solid green arrows represent toll-free travel, while dashed red arrows indicate that a toll is applied in the direction of the arc.

\begin{figure}[h!]
\centering
\includegraphics[scale = 0.4]{Figures/new-york-tolls}
\caption{Simplified NYC Transportation Network with Directional Tolls}
\label{fig:nyc_network}
\end{figure}

\subsection{Either Or Constraints}
\emph{``At least one of these constraints holds"} is what we would like to model.  Equivalently, we can phrase this as an \emph{inclusive or} constraint.  This can be molded with a pair of Big-M constraints.
\begin{general}{Either Or}{}{}

\begin{equation}
\text{Either} \ \ \ a^\top x \leq b\ \  \text{  or  } \ \  c^\top x \leq d \ \ \text{holds} 
\end{equation}
can be modeled as 
\begin{equation}
\begin{split}
a^\top x - b &\leq M_1 \delta\\
c^\top x - d &\leq M_2 (1-\delta)\\
\delta &\in \{0,1\},
\end{split}
\end{equation}
where $M_1$ is an upper bound on $a^\top x - b$ and $M_2$ is an upper bound on $c^\top x - d$.
\end{general}

\paragraph{Interpretation.}
\begin{itemize}
    \item If \(\delta = 0\), then the first constraint becomes \(a^\top x \leq b\), and the second becomes \(c^\top x \leq d + M_2\), which is trivially satisfied if \(M_2\) is large enough.
    \item If \(\delta = 1\), then the second constraint becomes \(c^\top x \leq d\), while the first is relaxed as \(a^\top x \leq b + M_1\).
\end{itemize}

\paragraph{Choosing Big-M Values.}
The constants \(M_1\) and \(M_2\) should be chosen as valid upper bounds on the expressions \(a^\top x - b\) and \(c^\top x - d\), respectively, over the feasible region. Overly large values may weaken the LP relaxation and cause numerical instability, so careful bounding is recommended whenever possible.

\paragraph{Use Cases.}
Either-or constraints naturally arise in problems with conditional decisions, such as:
\begin{itemize}
    \item If a machine is used, then either a certain energy constraint or a safety constraint must be met.
    \item In facility planning, either a site must be built with one configuration or another.
    \item In piecewise models where a variable must fall into one of two feasible regions.
\end{itemize}

\begin{example}{}{}
Either $2$ buses or $10$ cars are needed shuttle students to the football game.  
\begin{itemize}
\item Let $x$ be the number of buses we have and 
\item let $y$ be the number of cars that we have.  
\end{itemize}
Suppose that there are at most $M_1 = 5$ buses that could be rented an at most $M_2 = 20$ cars that could be available.

This constraint can be modeled as 
\begin{equation}
\begin{split}
x - 2 &\leq 5\delta\\
y - 10 &\leq 20 (1-\delta)\\
\delta &\in \{0,1\},
\end{split}
\end{equation}
\end{example}
\subsection{If then implications - opposite direction}
Suppose that we want to model the fact that if we have at most 10 students attending this course, then we must switch to a smaller classroom.   

Let $x_i \in \{0,1\}$ be 1 if student $i$ is in the course or not.   Let $\delta \in \{0,1\}$ be 1 if we need to switch to a smaller classroom.  

Thus, we want to model 
\begin{quote}
If $$
\sum_{i \in I} x_i \leq 10
$$
then 
$$
\delta = 1.
$$
\end{quote}

We can model this as 
\begin{equation}
\sum_{i \in I} x_i \geq 10 + 1 + M\delta.
\end{equation}

\begin{general}{If inequality, then indicator}{}
We let $m$ be a lower bound on the quantity $a^\top x - b$ and we let $\epsilon$ be a tiny number that is an error bound in verifying if an inequality is violated.  \textbf{If the data $a,b$ are integer and $x$ is an integer, then we can take $\epsilon = 1$.}

Now
\begin{equation}
\text{If } \ \ a^\top x \leq b  \ \ \text{then}\ \ \delta = 1
\end{equation}
can be modeled as 
\begin{equation}
a^\top x -b  \geq  \epsilon(1-\delta) + m \delta.
\end{equation}
\end{general}

\begin{proof}We not justify the statement above.

A simple way to understand this constraint is to consider the \emph{contrapositive} of the if then statement that we want to model.  The contrapositive says that 
\begin{equation}
\text{If $\delta = 0$, then $a^\top x - b > 0$.}
\end{equation}
To show the contrapositive, we set $\delta = 0$.  Then the inequality becomes 
$$
a^\top x - b \geq \epsilon(1-0) + m0 = \epsilon > 0.
$$
Thus, the contrapositive holds.

\textbf{If instead we wanted a direct proof:}

Case 1: Suppose $a^\top x \leq b$.  Then $0 \geq a^\top x - b$, which implies that 
$$
\delta(a^\top x - b) \geq a^\top x - b
$$
Therefore
$$
\delta(a^\top x - b) \geq \epsilon(1-\delta) + m \delta
$$
After rearranging
$$
\delta(a^\top x - b - m) \geq \epsilon(1-\delta)
$$
Since $a^\top x - b - m \geq 0$ and $\epsilon > 0$, the only feasible choice is $\delta = 1$.

Case 2:  Suppose $a^\top x > b$.  Then $a^\top x - b \geq \epsilon$.    Since $a^\top x -b \geq m$, both choices $\delta = 0$ and $\delta = 1$ are feasible.

By the choice of $\epsilon$,  we know that $a^\top x -b > 0$ implies that $a^\top x - b \geq \epsilon$.  

Since we don't like strict inequalities, we write the strict inequality as $a^\top x - b \geq \epsilon$ where $\epsilon$ is a small positive number that is a smallest difference between $a^\top x - b$ and $0$ that we would typically observe.  As mentioned above, if $a,b,x$ are all integer, then we can use $\epsilon = 1$.

Now we want an inequality with left hand side $a^\top x - b \geq$ and right hand side to take the value 
\begin{itemize}
\item $\epsilon$ if $\delta = 0$,
\item $m$ if $\delta = 1$.
\end{itemize}
This is accomplished with right hand side $\epsilon (1-\delta) + m\delta$.
\end{proof}

Many other combinations of if then statements are summarized in the following table:
\begin{table}
\begin{center}
\begin{tabular}{|c|c|}
\hline
\textbf{Implication} & \textbf{Constraint}\\
\hline
If $\delta = 0$, then $a^\top x \leq b$ & $a^\top x \leq b + M \delta$\\
If $a^\top x \leq b$, then $\delta = 1$ & $a^\top x \geq m \delta + \epsilon(1-\delta)$\\
\hline
\end{tabular}
\end{center}
\caption{Short list: If/then models with a constraint and a binary variable.  Here $M$ and $m$ are upper and lower bounds on $a^\top x - b$ and $\epsilon$ is a small number such that if $a^\top x > b$, then $a^\top x \geq b + \epsilon$.}
\end{table}
These two implications can be used to derive the following longer list of implications.

\begin{table}
\begin{center}
\begin{tabular}{|c|c|}
\hline
\textbf{Implication} & \textbf{Constraint}\\
\hline
If $\delta = 0$, then $a^\top x \leq b$ & $a^\top x \leq b + M \delta$\\
If $\delta = 0$, then $a^\top x \geq b$ & $a^\top x \geq b + m \delta$\\
If $\delta = 1$, then $a^\top x \leq b$ & $a^\top x \leq b + M (1-\delta)$\\
If $\delta = 1$, then $a^\top x \geq b$ & a$^\top x \geq b + m (1-\delta)$\\
If $a^\top x \leq b$, then $\delta = 1$ & $a^\top x \geq b + m \delta + \epsilon(1-\delta)$\\
If $a^\top x \geq b$, then $\delta = 1$ & $a^\top x \leq b + M \delta - \epsilon(1-\delta)$\\
If $a^\top x \leq b$, then $\delta = 0$ & $a^\top x \geq b + m (1-\delta) + \epsilon \delta$\\
If $a^\top x \geq b$, then $\delta = 0$ & $a^\top x \geq b + m (1-\delta) - \epsilon \delta$\\
\hline
\end{tabular}
\end{center}
\caption{Long list: If/then models with a constraint and a binary variable.  Here $M$ and $m$ are upper and lower bounds on $a^\top x - b$ and $\epsilon$ is a small number such that if $a^\top x > b$, then $a^\top x \geq b + \epsilon$.}
\end{table}
Lastly, if you insist on having exact correspondance, that is, "$\delta = 0 $ if and only if $a^\top x \leq b$" you can simply include both constraints for "if $\delta = 0 $, then $a^\top x \leq b$" and if "$a^\top x \leq b$, then $\delta = 0 $".  Although many problems may be phrased in a way that suggests you need "if and only if", it is often not necessary to use both constraints due to the objectives in the problem that naturally prevent one of these from happening.  

For example, if we want to add a binary variable $\delta$ that means
$$
\begin{cases}
\delta = 0 \text{ implies }  a^\top x \leq b\\
\delta = 1  \text{ Otherwise}
\end{cases}
$$
If $\delta = 1$ does not effect the rest of the optimization problem, then adding the constraint regarding $\delta = 1$ is not necessary.  Hence, typically, in this scenario, we only need to add the constraint $a^\top x \leq b + M \delta$.

\subsection{Multi Term Disjunction with application to 2D packing}

\section{Generalized $k$-of-$n$ Disjunction}
A disjunction is a generalization of an ``or'' statement. Suppose that we have $n$ constraints 
	$$
	\mathbf{a}_1^\top\mathbf{x} \leq b_1, \quad \mathbf{a}_2^\top\mathbf{x} \leq b_2, \quad\dots, \quad \mathbf{a}_n^\top\mathbf{x} \leq b_n
	$$
and we want to enforce at least $k$ of them. This can be accomplished linearly by introducing a new binary indicator variable $\delta_i$ for each of the disjunctive constraints $i \in \{1,\dots,n\}$:
	\begin{align*}
	\mathbf{a}_1^\top\mathbf{x} &\leq b_1+M(1-\delta_1) \\
	\mathbf{a}_2^\top\mathbf{x} &\leq b_2+M(1-\delta_2) \\
								&\vdots 				\\
	\mathbf{a}_n^\top\mathbf{x} &\leq b_n+M(1-\delta_1) \\
	\sum_{i=1}^n\delta_i &\geq k
	\end{align*}
If $\delta_i = 1$, then the $i^\text{th}$ disjunctive constraint is actively enforced. The last constraint $\left(\sum_{i=1}^n\delta_i \geq 1\right)$ ensures that at least $k$ of the constraints is active.

\subsubsection{Strip Packing Problem}
Suppose we a selection of rectangles $I$ and a 2-dimensional strip with width $W$ and infinite height. Each rectangle $i \in I$ has width $w_i$ and 
height $h_i$ and we want to pack the rectangles into the strip so that \eqref{mod:StripPackingObj-simple} overall height is minimized, \eqref{mod:StripPackingWidth-simple} overall width is less than $W$, and \eqref{mod:StripPackingOverlap-simple} none of the rectangles overlap. See Figure \ref{fig:StripPacking} for an example.
	\begin{figure}[h!]
	\begin{subfigure}[b]{.6\textwidth}
	\centering
	\includegraphics[height=.4\textheight]{Figures/StripPacking0}
	\caption{Pack the rectangles into the strip}
	\label{fig:StripPacking-PackStrip}
	\end{subfigure}
	\hfill
	\begin{subfigure}[b]{.3\textwidth}
	\centering
	\includegraphics[height=.4\textheight]{Figures/StripPacking1}
	\caption{Minimize overall height $H$.}
	\label{fig:StripPacking-MinHeight}
	\end{subfigure}
	\caption{An Example of the Strip Packing Problem}
	\label{fig:StripPacking}
	\end{figure}

Let $(x_i,y_i)$ denote the position of the lower-left-hand corner of each rectangle $i \in I$. The overlapping constraint \eqref{mod:StripPackingOverlap-simple} is the trickiest part. Consider a pair of rectangles $i$ and $j$: rectangle $j$ is located entirely to the left of rectangle $i$ if $x_j$ has a value larger than $x_j+w_j$. That is, if $x_j \geq x_j+w_j$. On the other hand, if $y_j \geq y_i+h_i$, then rectangle $j$ is located entirely above rectangle $i$. If either of these constraints is satisfied then rectangles $i$ and $j$ do not overlap. We could also place $i$ above or to the right of $j$. This gives four constraints that we need to satisfy at least one of. The model can be thought of thusly:

	\begin{subequations}
	\begin{alignat}{4}
	\text{Minimize}& \quad & \mathclap{\hspace{-69pt} \max_{i\in I}\{y_i+h_i\}} 	\label{mod:StripPackingObj-simple}\\
		\text{s.t.}& 	   & \max_{i\in I}\{x_i+w_i\} &\leq W 						\label{mod:StripPackingWidth-simple}\\
				   &	   & \mathclap{\hspace{111pt} \left.\begin{aligned}
				   		x_i+w_i &\leq x_j \\
				   		x_j+w_j &\leq x_i \\
				   		y_i+h_i &\leq y_j \\
				   		y_j+h_j &\leq y_i
				   		\end{aligned}\right\} \parbox{110pt}{At least one of these for every distinct pairs of rectangles $i,j \in I$}} \label{mod:StripPackingOverlap-simple} \\
			   	   &	   & x_i,y_i &\geq 0		&\qquad& & &\forall\ i\in I
	\end{alignat}
	\end{subequations}
Constraint \eqref{mod:StripPackingOverlap-simple} is a set of four \textit{disjunctive} constraints. This can be expressed linearly by introducing a new binary indicator variable $\delta_{ijk}$ for each of the disjunctive constraints in \eqref{mod:StripPackingOverlap-simple}:
	\begin{alignat}{4}
	\text{Minimize}& \quad & \mathclap{\hspace{-53pt} H}	\nonumber\\
		\text{s.t.}& 	   & y_i+h_i &\leq H						&\qquad& & &\forall\ i\in I		\tag{\ref{mod:StripPackingObj-simple}}	\\
				   & 	   & x_i+w_i &\leq W						&\qquad& & &\forall\ i\in I		\tag{\ref{mod:StripPackingWidth-simple}}\\
				   &	   & x_i+w_i &\leq x_j+M(1-\delta_{ij1}) 	&\qquad& & &\forall\ i,j\in I:\ i>j	\nonumber \\
				   &	   & x_j+w_j &\leq x_i+M(1-\delta_{ij2}) 	&\qquad& & &\forall\ i,j\in I:\ i>j	\nonumber \\
				   &	   & y_i+h_i &\leq y_j+M(1-\delta_{ij3}) 	&\qquad& & &\forall\ i,j\in I:\ i>j	\tag{\ref{mod:StripPackingOverlap-simple}}\\
				   &	   & y_j+h_j &\leq y_i+M(1-\delta_{ij4}) 	&\qquad& & &\forall\ i,j\in I:\ i>j	\nonumber \\
				   &	   & \sum_{k=1}^4\delta_{ijk} & \geq 1	 	&\qquad& & &\forall\ i,j\in I:\ i>j	\nonumber \\
				   &	   & x_i,y_i &\geq 0		&\qquad& & &\forall\ i\in I							\nonumber \\
				   &	   & \boldsymbol{\delta}_{ij} &\in\{0,1\}^4	&\qquad& & &\forall\ i,j\in I:\ i>j	\nonumber 
	\end{alignat}
If $\delta_{ijk} = 1$, then the $k^\text{th}$ disjunctive constraint is actively enforced. The last constraint $\left(\sum_{k=1}^4\delta_{ijk} \geq 1\right)$ ensures that at least one of the constraints is active. An optimal solution to the example is given in Figure \ref{fig:StripPacking-Opt}; it was found via GurobiPy.
	\begin{figure}[h!]
	\centering
	\includegraphics[height=.4\textheight]{Figures/StripPacking2}
	\caption{The Optimal Solution to the example problem}
	\label{fig:StripPacking-Opt}
	\end{figure}

This problem is considered strongly NP-Hard.

\subsection{SOS1 Constraints}
\begin{definition}{Special Ordered Sets of Type 1 (SOS1)}
A Special Ordered Sets of type 1 (SOS1) constraint on a vector indicates that \emph{at most one element of the vector can non-zero}.
\end{definition}
We next give an example of how to use binary variables to model this and then show how much simpler it can be coded using the SOS1 constraint.

\begin{examplewithcode}{SOS1 Constraints}{\href{https://github.com/open-optimization/open-optimization-or-examples/blob/master/nonlinear-programming/Gurobi-SOS-Examples.ipynb}{Gurobipy}}
\label{example:sos1}
Solve the following optimization problem: \[\begin{aligned}
\text{maximize}\quad & 3x_1 + 4x_2 + x_3 + 5x_4 \\
\text{subject to}\quad & 0 \le x_i \le 5 \\
& \text{at most one of the $x_i$ can be nonzero}
\end{aligned}\]
\end{examplewithcode}
\subsection{SOS2 Constraints}
\begin{definition}{Special Ordered Sets of Type 2 (SOS2)}{}
A \emph{Special Ordered Set of Type 2 (SOS2)} constraint on a vector indicates that \emph{at most two elements of the vector can non-zero AND the non-zero elements must appear consecutively}.
\end{definition}
We next give an example of how to use binary variables to model this and then show how much simpler it can be coded using the SOS2 constraint.
\begin{examplewithcode}{SOS2}{\href{https://github.com/open-optimization/open-optimization-or-examples/blob/master/nonlinear-programming/Gurobi-SOS-Examples.ipynb}{Gurobipy}}
\label{example:SOS2}
Solve the following optimization problem: \[\begin{aligned}
\text{maximize}\quad & 3x_1 + 4x_2 + x_3 + 5x_4 \\
\text{subject to}\quad & 0 \le x_i \le 5 \\
& \text{at most two of the $x_i$ can be nonzero} \\
& \text{and the nonzero $x_i$ must be consecutive}
\end{aligned}\]
\end{examplewithcode}

\subsection{Piecewise linear functions with SOS2 constraint}
\begin{examplewithcode}{Piecewise Linear Function}{\href{https://github.com/open-optimization/open-optimization-or-examples/blob/master/nonlinear-programming/Gurobi-SOS-Examples.ipynb}{Gurobipy}}
\label{example:pwl}
Consider the piecewise linear function 
 $c(x)$ given by
$$
c(x) = 
\begin{cases}
25x  & \text{ if } 0 \leq x \leq 5\\
20x + 25 & \text{ if } 5 \leq x \leq 10\\
15x + 75 & \text{ if } 10 \leq x \leq 15
\end{cases}
$$
\begin{center}
\includegraphics[scale = 1]{Figures/pwl-plot.pdf}
\end{center}
%See \refincludefigurestatic[scale = 0.35]{pwl-plot.png}.

We will use integer programming to describe this function.  We will fix $x = a$ and then the integer program will set the value $y$ to $c(a)$.
\begin{align*}
\min \quad & 0 \\
\st        & x = 5 z_{2} +10 z_{3} + 15 z_{4} \\
 & y = 125 z_{2} + 225 z_{3} + 300 z_{4} \\
 & z_{1} + z_{2} + z_{3} + z_{4} = 1 \\
 & SOS2: \{z_1, z_2, z_3, z_4\} \\
 & 0 \leq z_{i} \leq 1 &&\quad \text{for all } i \in \{1,2,3,4\} \\
 & x = a \\
\end{align*}

\end{examplewithcode}

\begin{examplewithcode}{Piecewise Linear Function Application}{code:pwl-application}
\label{example:pwl-application}
Consider the following optimization problem where the objective function includes the term $c(x)$, where $c(x)$ is the piecewise linear function described in \nameref{example:pwl}:
\begin{align}
\max \quad & z = 12x_{11} + 12x_{21} + 14x_{12} + 14x_{22} - c(x) \\
\st        & x_{11} + x_{12} \leq x + 5 \\
& x_{21} + x_{22} \leq 10 \\
& 0.5 x_{11} - 0.5x_{21} \geq 0 \\
& 0.4 x_{12} - 0.6 x_{22} \geq 0 \\
& x_{ij} \geq 0 \\
& 0 \leq x \leq 15
\end{align}

Given the piecewise linear, we can model the whole problem explicitly as a mixed-integer linear program.

 \begin{equation}
 \begin{array}{rrlr}
 \max\quad & 12 X_{1,1} + 12 X_{2,1} + 14 X_{1,2} + 14 X_{2,2} - y\\
\st \quad & x - 5 z_{2} - 10 z_{3} - 15 z_{4} &= 0\\
 & y - 125 z_{2} - 225 z_{3} - 300 z_{4} &= 0\\
 & z_{1} + z_{2} + z_{3} + z_{4} &= 1\\
 & X_{1,1} + X_{1,2} - x &\leq 5\\
 & X_{2,1} + X_{2,2} \leq 10\\
 & 0.5 X_{1,1} - 0.5 X_{2,1} &\geq 0\\
 & 0.4 X_{1,2} - 0.6 X_{2,2} &\geq 0\\
 & SOS2: \{z_1, z_2, z_3, z_4\}\\
 & X_{i,j} &\geq 0 &\quad\forall i \in \{1,2\}, j \in \{1,2\}\\
 & 0 \leq z_{i} &\leq 1 &\quad\forall i \in \{1,2,3,4\}\\
 & 0 \leq x &\leq 15\\
 & y & \text{ free} \\
\end{array}
\end{equation}
\end{examplewithcode}

\subsubsection{SOS2 with binary variables}

{Modeling Piecewise linear function}
{\begin{itemize}
\item Write down pairs of breakpoints and functions values $(a_i, f(a_i))$.
\item Define a binary variable $z_i$ indicating if $x$ is in the interval $[a_i, a_{i+1}]$.
\item Define multipliers $\lambda_i$ such that $x$ is a combination of the $a_i$'s and therefore the output $y = f(x)$ is a combination of the $f(a_i)$'s.
\item Restrict that at most 2 $\lambda_i's$ are non-zero and that those 2 are consecutive.

\end{itemize}

\begin{align*}
\min \quad & \sum_{i=1}^k \lambda_i f(a_i) \\
\st        & \sum_{i=1}^k \lambda_i = 1 \\
& x = \sum_{i=1}^k \lambda_i a_i \\
& \lambda_1 \leq z_1 \\
& \lambda_i \leq z_{i-1} + z_{i} &&\quad \text{for all } i=2, \dots, k-1, \\
& \lambda_k \leq z_{k-1} \\
& \lambda_i \geq 0, y_i \in \{0,1\}.
\end{align*}

\subsection{Maximizing a minimum}
When the constraints could be general, we will write $x \in X$ to define general constraints.  For instance, we could have $X = \{ x \in \R^n : Ax \leq b\}$ of $X  = \{ x \in \R^n : Ax \leq b, x \in \Z^n\}$ or many other possibilities.  

Consider the problem 

\begin{align*}
\max   \quad & \min \{x_1, \dots, x_n\}\\
\text{ such that } \quad &  x \in X
\end{align*}
Having the minimum on the inside is inconvenient.  To remove this, we just define a new variable $y$ and enforce that $y \leq x_i$ and then we maximize $y$.  Since we are maximizing $y$, it will take the value of the smallest $x_i$.  Thus, we can recast the problem as

\begin{align*}
\max\quad    & y\\
\text{ such that } \quad  & y \leq x_i \ \text{ for }\  i=1, \dots, n \\
&  x \in X
\end{align*}

\subsection{Relaxing (nonlinear) equality constraints}

There are a number of scenarios where the constraints can be relaxed without sacrificing optimal solutions to your problem.   In a similar vein of the maximizing a minimum, if because of the objective we know that certain constraints will be tight at optimal solutions, we can relax the equality to an inequality.   For example, 
\begin{align*}
\max   \quad &x_1 + x_2 +  \dots + x_n\\
\text{ such that } \quad &  x_i = y_i^2 + z_i^2 \text{ for } i=1, \dots, n
\end{align*}

%
\subsection{Connecting to continuous variables}
%
%Let $x_i \geq 0$ and $y_i \in \{0,1\}$ for all $i=1, \dots, n$.
%
\begin{enumerate}
\item If $x_i > 0$, then $y_i = 1$.
\begin{equation}
%x_i \leq M y_i
\end{equation}
%where $M$ is a suffiently large upper bound on the variable $x_i$.
\item If $x_i = 0$, then $y_i = 0$.\\
%This is harder to model!  Alternatively, we try modeling "if $x_i$ is sufficiently small, then $y_i = 0$.   For instance, if $x_i \leq 0.0000001$, then $y_i = 0$.
%This can be modeled as
\begin{equation}
%x_i - 0.0000001 \geq y_i -1.
\end{equation}
\item If $y_i = 1$, then $x_i \geq 5$
\begin{equation}
%5y_i \leq x_i.
\end{equation}
\end{enumerate}

\subsection{Modeling the Exact Absolute Value}

In many optimization models, we need to enforce an exact equality of the form $|x| = t$, where $x$ is a decision variable and $t \geq 0$.  Unlike the relaxation $|x| \leq t$ (which is convex and can be handled with two linear inequalities), the exact equality $|x| = t$ describes a \emph{non-convex} set and cannot be represented using linear constraints alone.

To model this, we decompose $x$ into its positive and negative parts: let $x^+, x^- \geq 0$ with
\[
x = x^+ - x^-, \qquad t = x^+ + x^-.
\]
This decomposition correctly gives $|x| = t$ provided that $x^+$ and $x^-$ are not both strictly positive simultaneously.  We enforce this complementarity condition using a binary variable $z \in \{0,1\}$ and a large constant $M$ (an a priori upper bound on $|x|$):
\[
\begin{aligned}
x &= x^+ - x^-, \\
t &= x^+ + x^-, \\
x^+ &\leq Mz, \\
x^- &\leq M(1 - z), \\
x^+, x^- &\geq 0, \quad z \in \{0,1\}.
\end{aligned}
\]
When $z = 1$, the constraint $x^- \leq 0$ forces $x^- = 0$, so $x = x^+ \geq 0$.  When $z = 0$, the constraint $x^+ \leq 0$ forces $x^+ = 0$, so $x = -x^- \leq 0$.  In either case, exactly one of $x^+$ and $x^-$ is active, and $t = |x|$ holds exactly.

\subsubsection{Modeling the Exact 1-Norm}

Building on the absolute value formulation, we can model the exact $\ell_1$-norm equality
\[
\sum_{i=1}^n |x_i| = c,
\]
where $\mathbf{x} \in \mathbb{R}^n$ is a vector of decision variables and $c > 0$ is a constant.  Constraints of this type arise, for instance, in portfolio optimization models that require full investment (where short-selling is permitted).

We introduce positive and negative parts for each component: $x_i = x_i^+ - x_i^-$ with $x_i^+, x_i^- \geq 0$.  Binary variables $z_i \in \{0,1\}$ enforce that at most one of $x_i^+$ and $x_i^-$ is nonzero:
\[
\begin{aligned}
x_i &= x_i^+ - x_i^-, \quad i = 1, \ldots, n, \\
x_i^+ &\leq c \, z_i, \quad i = 1, \ldots, n, \\
x_i^- &\leq c \, (1 - z_i), \quad i = 1, \ldots, n, \\
\sum_{i=1}^n (x_i^+ + x_i^-) &= c, \\
x_i^+, x_i^- &\geq 0, \quad z_i \in \{0,1\}, \quad i = 1, \ldots, n.
\end{aligned}
\]
Here we use $c$ as the big-$M$ bound, since no individual component can exceed $c$ in absolute value when the total 1-norm equals $c$.

\subsubsection{Modeling the Exact Maximum}

Suppose we wish to model the exact equality $t = \max\{x_1, \ldots, x_n\}$, where $x_1, \ldots, x_n$ and $t$ are decision variables.  The relaxed version $t \geq x_i$ for all $i$ only gives $t \geq \max_i x_i$; to enforce equality, we need to ensure that $t$ actually equals one of the $x_i$.

We introduce binary indicator variables $z_1, \ldots, z_n$, where $z_i = 1$ means that $x_i$ achieves the maximum.  Choosing a sufficiently large constant $M$, the formulation is:
\[
\begin{aligned}
x_i &\leq t, \quad i = 1, \ldots, n, \\
t &\leq x_i + M(1 - z_i), \quad i = 1, \ldots, n, \\
z_1 + z_2 + \cdots + z_n &= 1, \\
z_i &\in \{0,1\}, \quad i = 1, \ldots, n.
\end{aligned}
\]
The first set of constraints ensures $t \geq x_i$ for all $i$.  The second set, combined with exactly one $z_i = 1$, ensures that for the index $i$ with $z_i = 1$, we have $t \leq x_i$, which together with $t \geq x_i$ gives $t = x_i$.  For all other indices $j \neq i$, the big-$M$ term relaxes the upper bound constraint.

\newpage

\section{Job Shop Scheduling}

Example borrowed from: \href{https://python-mip.readthedocs.io/en/latest/examples.html#job-shop-scheduling-problem}{Pyhton MIP example}

The Job Shop Scheduling Problem, abbreviated as JSSP, is a classical combinatorial optimization problem that is classified as NP-hard. It is characterized by a collection of jobs that are to be processed on a specific group of machines. Each job has a predetermined processing sequence across the machines, with a fixed processing duration for each machine. A job can only visit a machine once, and machines can only accommodate one job at a time. Once a machine starts a job, it must finish it without any interruptions. The primary goal is to reduce the makespan, which is the total time taken to complete all jobs.

Consider the case of having three jobs and three machines. The jobs have the following processing sequences and times:

\begin{itemize}
  \item Job \( j_1 \): \( m_3(2) \) followed by \( m_1(1) \) and then \( m_2(2) \).
  \item Job \( j_2 \): \( m_2(1) \) followed by \( m_3(2) \) and then \( m_1(2) \).
  \item Job \( j_3 \): \( m_3(1) \) followed by \( m_2(2) \) and then \( m_1(1) \).
\end{itemize}

Below, we present two potential schedules. The first one is a straightforward approach where jobs are processed sequentially, leading to machines having idle periods. The second approach is optimal, where jobs are executed concurrently.

\subsection{JSSP Components}
The JSSP can be formally described by the following parameters:

\begin{itemize}
  \item \( J\): Represents the set of jobs, indexed as \( \{1, \ldots, n\} \).
  \item \( I\): Represents the set of machines, indexed as \( \{1, \ldots, m\} \).
  \item \( o_r^j \): Specifies the machine that handles the \( r \)-th operation of job \( j \). The sequence \( O^j=\left(o_1^j, o_2^j, \ldots, o_m^j\right) \) illustrates the processing order for job \( j \).
  \item \( p_{i j} \): Represents the non-negative processing time of job \( j \) on machine \( i \).
\end{itemize}

The solution to the JSSP must satisfy:

\begin{enumerate}
  \item Every job \( j \) should be processed according to the machine sequence \( O^j \).
  \item At any given time, a machine can only process one job.
  \item A job, once started on a machine, must run to completion without any breaks.
\end{enumerate}

The objective is to minimize the makespan, or the finishing time of the last job. The problem remains NP-hard even if \( n \geq 3 \) or \( m \geq 3 \).

Decision variables include:
\begin{itemize}
  \item \( x_{i j} \): Starting time of job \( j \) on machine \( i \).
  \item \( y_{i j k} \): A binary variable. It's 1 if job \( j \) precedes job \( k \) on machine \( i \) and 0 otherwise.
  \item \( C \): The makespan or the maximum completion time.
\end{itemize}

\textbf{Example: Job Shop Scheduling Problem Input Data}

\textbf{Basic Parameters}
The number of jobs, \( n \), is 5 and the number of machines, \( m \), is 3. 

\textbf{Processing Times}
The matrix of processing times for each job on each machine is given by:

\[
\text{times} = 
\begin{bmatrix}
2 & 1 & 2 \\
1 & 2 & 2 \\
1 & 2 & 1 \\
1 & 1 & 1 \\
3 & 4 & 5 \\
\end{bmatrix}
\]

Where the element in the \(i\)-th row and \(j\)-th column represents the time required to process job \(i\) on machine \(j\).

\textbf{Machine Sequences}
The order in which each job needs to be processed on the machines is described by:

\[
\text{machines} = 
\begin{bmatrix}
2 & 0 & 1 \\
1 & 2 & 0 \\
2 & 1 & 0 \\
1 & 2 & 0 \\
0 & 1 & 2 \\
\end{bmatrix}
\]

Here, the sequence of machines for each job is given by the rows of the matrix. For instance, job 1 is processed first on machine 2, then on machine 0, and finally on machine 1.

\textbf{Big \(M\) Calculation}
A large constant, \( M \), is used in the model formulation. It is computed as the sum of all processing times, given by:

\[
M = \sum_{i=1}^{n} \sum_{j=1}^{m} \text{times}[i][j]
\]

\begin{figure}
% Will use file.tikz if it exists, else file.pdf
\includetikz[width=0.8\textwidth]{Figures/jssp}
\end{figure}
\refincludefigurestatic[scale = 1.5]{jssp.png}

\paragraph{Applications of JSSP}
The JSSP has numerous practical applications, including:
\begin{itemize}
  \item Manufacturing: Scheduling tasks in a production line to optimize throughput.
  \item Computer Science: Allocating tasks to processors in parallel computing environments.
  \item Healthcare: Scheduling surgeries and other treatments in hospitals.
  \item Transportation: Optimizing maintenance tasks for vehicles or aircraft in a maintenance facility.
\end{itemize}

\subsection{Mathematical Model}

\modelhead{Sets}
\begin{itemize}
  \item \( J\): Represents the set of jobs, indexed as \( \{1, \ldots, n\} \).
  \item \( I \): Represents the set of machines, indexed as \( \{1, \ldots, m\} \).
  \end{itemize}
\modelhead{Parameters}
\begin{itemize}
 \item \( o_r^j \): Specifies the machine that handles the \( r \)-th operation of job \( j \). The sequence \( O^j=\left(o_1^j, o_2^j, \ldots, o_m^j\right) \) illustrates the processing order for job \( j \).
  \item \( p_{i j} \): Represents the non-negative processing time of job \( j \) on machine \( i \).
\end{itemize}
\modelhead{Variables}
$$
x_{ij} = \text{ start time of job $j$ on machine $i$}.
$$

$$
y_{ijk} = \begin{cases} 1, & \text{if job } j \text{ precedes job } k \text{ on machine } i \text{,}\\ & i \in I \text{, } j, k  \in J \text{, } j \neq k \\ 0, & \text{otherwise} \end{cases}
$$

\modelhead{Model}
\begin{subequations}
\begin{equation}
\begin{array}{rrlrr}
\min 
               & C_{\max} \\
\st \\
&& \text{Overlapping constraints}\\
&            x_{o^{j}_{r}j} &  \geq x_{o^{j}_{r-1}j} +p_{o^{j}_{r-1}j} & \forall r \in \{2,..,m\}, j \in {J} & \text{Order tasks for job $j$} \\
             &      x_{ij}     & \geq x_{ik} + p_{ik} - M \cdot y_{ijk} & \forall j,k \in {J}, j \neq k, i \in {I} & \text{Ordering on machine $i$} \\
              &     x_{ik}     & \geq x_{ij} + p_{ij} - M \cdot (1-y_{ijk}) & \forall j,k \in {J}, j \neq k,i \in {I} & \text{Reverse ordering on machine $i$}\\
              \\
           &   & \text{ Bounding objective}\\
            &       C_{\max}          & \geq x_{o^{j}_{m}j} + p_{o^{j}_{m}j} & \forall j \in {J} & \text{last task of job $j$ completed before end}\\
            \\
            && \text{ Domain of variables}\\
            &      x_{ij}      & \geq 0 & \forall i \in {J}, i \in {I} \\
           &       y_{ijk}     & \in \{0,1\} & \forall j,k \in {J}, i \in {I} \\
          &        C_{\max} & \geq 0
              \end{array}
              \end{equation}
\end{subequations}
\begin{examplewithallcode}{Duplo Scheduling}{}{}{https://github.com/open-optimization/open-optimization-or-examples/blob/master/integer-programming/job-shop-scheduling.ipynb}

The Duplo in class exercise is solved with Gurobi in the link.

The optimal value is 11.

\begin{center}
\includetikz[width = 0.8\textwidth]{Figures/jssp-duplo}
\includegraphics[scale = 0.5]{optimization/figures/figures-static/jssp-duplo}

\includegraphics[scale=0.07, angle=180]{optimization/figures/figures-static/jssp-duplo-actual}
\end{center}
\end{examplewithallcode}

\subsection{Job Shop Scheduling Variations}

Makespan minimization of assining jobs to machines.

In this variation, each machine can handle a number of of different types of jobs.  Some machines can do certain jobs faster than others. 

We want to minimize the completion time of all of the jobs.  

In this simple variation, each job only needs to be processed on a single machine, but that choice of machine might vary dependent on which machine can do that job faster and which one is available.

%
%
%Let there be $n$ jobs and $m$ machines.
%
%Let $x_{ij}$ be a binary variable that is 1 if job $i$ is assigned to machine $j$, and 0 otherwise.
%
%Subject to:
%
%Each job can be assigned to only one machine: $\sum\limits_{j=1}^{m} x_{ij} = 1$ for $i=1$ to $n$
%
%Each machine can process only one job at a time: $\sum\limits_{i=1}^{n} x_{ij} \leq 1$ for $j=1$ to $m$
%
%Start time of a job is based on the completion time of its predecessor: $y_i \geq y_j + p_j$ for $(i,j) \in E$
%
%$y_i \geq 0$ for $i=1$ to $n$
%
%$x_{ij} \in {0,1}$ for $i=1$ to $n$ and $j=1$ to $m$
%
%$y_i \in \mathbb{N}$ for $i=1$ to $n$
%
%Where, $p_i$ is the processing time of job $i$, $E$ is the set of edges that represent the precedence constraints between the jobs.
%
%The model can be implemented using python with Gurobi library as follows:
%
\begin{verbatim}
%from gurobipy import *
%
%def job_shop_scheduling(n,m,p,E):
%    # Create a new model
%    model = Model()
%
%    # Create variables
%    x = {}
%    y = {}
%    for i in range(n):
%        y[i] = model.addVar(vtype=GRB.INTEGER, name="y_"+str(i))
%        for j in range(m):
%            x[i,j] = model.addVar(vtype=GRB.BINARY, name="x_"+str(i)+"_"+str(j))
%
%    # Set objective
%    obj = quicksum([y[i]+p[i] for i in range(n)])
%    model.setObjective(obj,GRB.MINIMIZE)
%
%    # Add constraints
%    for i in range(n):
%        model.addConstr(quicksum([x[i,j] for j in range(m)]) == 1)
%
%    for j in range(m):
%        model.addConstr(quicksum([x[i,j] for i in range(n)]) <= 1)
%
%    for (i,j) in E:
%        model.addConstr(y[i] >= y[j]+p[j])
%
%    for i in range(n):
%        model.addConstr(y[i] >= 0)
%
%    # Optimize model
%    model.optimize()
%
%    return
\end{verbatim}

%Here is an example:
%Here is an example job shop scheduling problem with 5 machines and 15 jobs:
%
%Machine 1: Can perform jobs 1, 2, 3, 4, and 5
%Machine 2: Can perform jobs 6, 7, 8, 9, and 10
%Machine 3: Can perform jobs 11, 12, 13, 14, and 15
%Machine 4: Can perform jobs 1, 6, 11, 2, and 7
%Machine 5: Can perform jobs 3, 8, 13, 4, and 9
%Jobs:
%
%Job 1: Machine 1 (processing time = 2 hours), Machine 4 (processing time = 1 hour)
%Job 2: Machine 1 (processing time = 3 hours), Machine 4 (processing time = 2 hours)
%Job 3: Machine 1 (processing time = 2 hours), Machine 5 (processing time = 1 hour)
%Job 4: Machine 1 (processing time = 4 hours), Machine 5 (processing time = 2 hours)
%Job 5: Machine 1 (processing time = 1 hour)
%Job 6: Machine 2 (processing time = 3 hours), Machine 4 (processing time = 2 hours)
%Job 7: Machine 2 (processing time = 2 hours), Machine 4 (processing time = 1 hour)
%Job 8: Machine 2 (processing time = 4 hours), Machine 5 (processing time = 2 hours)
%Job 9: Machine 2 (processing time = 1 hour), Machine 5 (processing time = 1 hour)
%Job 10: Machine 2 (processing time = 2 hours)
%Job 11: Machine 3 (processing time = 3 hours), Machine 4 (processing time = 2 hours)
%Job 12: Machine 3 (processing time = 2 hours), Machine 4 (processing time = 1 hour)
%Job 13: Machine 3 (processing time = 4 hours), Machine 5 (processing time = 2 hours)
%Job 14: Machine 3 (processing time = 1 hour), Machine 5 (processing time = 1 hour)
%Job 15: Machine 3 (processing time = 2 hours)
%Objective: Minimize the makespan (i.e. the total time it takes to complete all the jobs)
%
%Constraints:
%
%Each job can only be performed on the specified machines in the order listed
%A machine can only work on one job at a time
%There are no precedence constraints between jobs.
%

\textbf{Machines:}
\begin{itemize}
    \item Machine 1: Can perform jobs 1, 2, 3, 4, and 5
    \item Machine 2: Can perform jobs 6, 7, 8, 9, and 10
    \item Machine 3: Can perform jobs 11, 12, 13, 14, and 15
    \item Machine 4: Can perform jobs 1, 6, 11, 2, and 7
    \item Machine 5: Can perform jobs 3, 8, 13, 4, and 9
\end{itemize}

\textbf{Jobs:}
\begin{itemize}
    \item Job 1: Machine 1 (processing time = 2 hours), Machine 4 (processing time = 1 hour)
    \item Job 2: Machine 1 (processing time = 3 hours), Machine 4 (processing time = 2 hours)
    \item Job 3: Machine 1 (processing time = 2 hours), Machine 5 (processing time = 1 hour)
    \item Job 4: Machine 1 (processing time = 4 hours), Machine 5 (processing time = 2 hours)
    \item Job 5: Machine 1 (processing time = 1 hour)
    \item Job 6: Machine 2 (processing time = 3 hours), Machine 4 (processing time = 2 hours)
    \item Job 7: Machine 2 (processing time = 2 hours), Machine 4 (processing time = 1 hour)
    \item Job 8: Machine 2 (processing time = 4 hours), Machine 5 (processing time = 2 hours)
    \item Job 9: Machine 2 (processing time = 1 hour), Machine 5 (processing time = 1 hour)
    \item Job 10: Machine 2 (processing time = 2 hours)
    \item Job 11: Machine 3 (processing time = 3 hours), Machine 4 (processing time = 2 hours)
    \item Job 12: Machine 3 (processing time = 2 hours), Machine 4 (processing time = 1 hour)
    \item Job 13: Machine 3 (processing time = 4 hours), Machine 5 (processing time = 2 hours)
    \item Job 14: Machine 3 (processing time = 1 hour), Machine 5 (processing time = 1 hour)
    \item Job 15: Machine 3 (processing time = 2 hours)
\end{itemize}

\textbf{Objective:} Minimize the makespan (i.e. the total time it takes to complete all the jobs)

\modelhead{Constraints}
\begin{itemize}
    \item Each job can only be performed on the specified machines in the order listed
    \item A machine can only work on one job at a time
    \item There are no precedence constraints between jobs.
\end{itemize}
%
\begin{verbatim}
%job_shop_data = {
%    'machines': {
%        1: {'jobs': [1, 2, 3, 4, 5]},
%        2: {'jobs': [6, 7, 8, 9, 10]},
%        3: {'jobs': [11, 12, 13, 14, 15]},
%        4: {'jobs': [1, 6, 11, 2, 7]},
%        5: {'jobs': [3, 8, 13, 4, 9]}
%    },
%    'jobs': {
%        1: {'machine_times': {1: 2, 4: 1}},
%        2: {'machine_times': {1: 3, 4: 2}},
%        3: {'machine_times': {1: 2, 5: 1}},
%        4: {'machine_times': {1: 4, 5: 2}},
%        5: {'machine_times': {1: 1}},
%        6: {'machine_times': {2: 3, 4: 2}},
%        7: {'machine_times': {2: 2, 4: 1}},
%        8: {'machine_times': {2: 4, 5: 2}},
%        9: {'machine_times': {2: 1, 5: 1}},
%        10: {'machine_times': {2: 2}},
%        11: {'machine_times': {3: 3, 4: 2}},
%        12: {'machine_times': {3: 2, 4: 1}},
%        13: {'machine_times': {3: 4, 5: 2}},
%        14: {'machine_times': {3: 1, 5: 1}},
%        15: {'machine_times': {3: 2}}
%    },
%    'objective': 'minimize_makespan',
%    'constraints': ['job_machine_mapping', 'single_job_single_machine', 'no_job_precedence']
%}
\end{verbatim}

We leave this as an exercise for the reader to model and solve.

\section{Quadratic Assignment Problem (QAP)}
\begin{resource}{}{}
\begin{itemize}
\item \href{https://ieeexplore-ieee-org.ezproxy.lib.vt.edu/stamp/stamp.jsp?tp=&arnumber=7170278}{An applied case of quadratic assignment problem in hospital department layout}
\item See \href{https://www.semanticscholar.org/paper/Quadratic-Assignment-Problem%3A-A-survey-and-Shawky-Metwally/247b45613d6d3b6961fdad44f9e7fefb70fd3e82}{Quadratic Assignment Problem: A survey and Applications}.
\item \href{https://vtechworks.lib.vt.edu/server/api/core/bitstreams/06b13a4f-d1a9-452c-957f-b0bdd1619c76/content}{Quadratic Assignment Problem in Airline Gate Assignment to minimize walking distance}
\end{itemize}
\end{resource}
The quadratic assignment problem must choose the assignment of $n$ facilities to $n$ locations.  Each facility sends some flow to each other facility, and there is a distance to consider between locations.   
The objective is to minimize to distance times the flow of the assignment.

\textbf{Example:  Hospital Layout}
On any given day in the hospital, there will be patients that move from various locations in the hospital to various other locations in the hospital.  For example, patients move from the operating room to a recovery room, or from the emergency room to the operating room, etc.

We would like to chose the locations of these places in the hospital to minimize the amount of total distance traveled by all the patients. 

\begin{general}{Quadratic Assignment Problem}{\npcomplete}
Given flow $f_{ij}$  connections $c_{ij}$ and fixed building costs $f_i$, demands $d_j$ and capacities $u_i$, the capacitated facility location problem is 

\modelhead{Sets}
\begin{itemize}
\item Let $I = \{1,\dots, n\}$ be the set of facilities.
\item Let $K = \{1, \dots, n\}$ be the set of locations.
\end{itemize}

\modelhead{Parameters}
\begin{itemize}
\item $f_{ij}$ - flow from facility $i$ to facility $j$.
\item $d_{kl}$ - distance from location $k$ to location $l$.
\item $c_{ik}$ - cost to setup facility $i$ at location $k$.
\end{itemize}

\modelhead{Variables}
\begin{itemize}
\item Let 
\begin{equation*}
x_{ik} = \begin{cases}
1 & \text{if we place facility $i$ in location $k$,}\\
0 & \text{otherwise.}
\end{cases}
\end{equation*}
\end{itemize}

\modelhead{Model}
\begin{align}
\min \quad & \displaystyle\sum_{i=1}^n\sum_{j=1}^n \sum_{k=1}^n \sum_{l = 1}^n f_{ij}d_{kl}x_{ik}x_{jl} + \sum_{i=1}^n \sum_{k = 1}^n c_{ik} x_{ik}   \tag{total cost} \\
\st        & \displaystyle\sum_{i=1}^n x_{ik}=1 &&\quad \text{for all } k=1,\dots,n   \tag{assign facility to location $k$} \\
& \displaystyle \sum_{k=1}^n x_{ik}=1 &&\quad \text{for all } i=1\dots,n   \tag{assign one location to facility $i$} \\
& x_{ik}\in\{0,1\} &&\quad \text{for all } i=1,\dots,n, \text{ and } k = 1, \dots, n   \tag{binary decisions}
\end{align}
\end{general}

\section{Genaralized Assignment Problem (GAP)}

\url{https://en.wikipedia.org/wiki/Generalized_assignment_problem}
In \href{applied_mathematics}{applied mathematics}, the maximum
\textbf{generalized assignment problem} is a problem in
\href{combinatorial_optimization}{combinatorial optimization}. This
problem is a \url{generalization} of the
\href{assignment_problem}{assignment problem} in which both tasks and
\href{Agent-based_model}{agents} have a size. Moreover, the size of each
task might vary from one agent to the other.

This problem in its most general form is as follows: There are a number
of agents and a number of tasks. Any agent can be assigned to perform
any task, incurring some cost and profit that may vary depending on the
agent-task assignment. Moreover, each agent has a budget and the sum of
the costs of tasks assigned to it cannot exceed this budget. It is
required to find an assignment in which all agents do not exceed their
budget and total profit of the assignment is maximized.

\hypertarget{in-special-cases}{%
\subsection{In special cases}\label{in-special-cases}}

In the special case in which all the agents' budgets and all tasks'
costs are equal to 1, this problem reduces to the
\href{assignment_problem}{assignment problem}. When the costs and
profits of all tasks do not vary between different agents, this problem
reduces to the multiple knapsack problem. If there is a single agent,
then, this problem reduces to the \href{knapsack_problem}{knapsack
problem}.

\hypertarget{explanation-of-definition}{%
\subsection{Explanation of definition}\label{explanation-of-definition}}

In the following, we have \emph{n} kinds of items, \(a_1\) through
\(a_n\) and \emph{m} kinds of bins \(b_1\) through \(b_m\). Each bin
\(b_i\) is associated with a budget \(t_i\). For a bin \(b_i\), each
item \(a_j\) has a profit \(p_{ij}\) and a weight \(w_{ij}\). A solution
is an assignment from items to bins. A feasible solution is a solution
in which for each bin \(b_i\) the total weight of assigned items is at
most \(t_i\). The solution's profit is the sum of profits for each
item-bin assignment. The goal is to find a maximum profit feasible
solution.

Mathematically the generalized assignment problem can be formulated as
an \href{Integer_programming}{integer program}:
\begin{align} \text{maximize } &
\sum_{i=1}^m\sum_{j=1}^n
p_{ij} x_{ij}. \\
\text{subject to } & \sum_{j=1}^n
w_{ij} x_{ij} \le t_i & & i=1,
\ldots, m; \\ &
\sum_{i=1}^m x_{ij} = 1 & & j=1,
\ldots, n; \\ & x_{ij}
\in \{0,1\} & & i=1,
\ldots, m, \quad j=1,
\ldots, n; \end{align}

\section{Other material}

\subsection{Binary reformulation of integer variables}
If an integer variable has small upper and lower bounds, it can sometimes be advantageous to recast it as a sequence of binary variables - for either modeling, the solver, or both.   Although there are technically many ways to do this, here are the two most common ways.

\begin{general}{Full reformulation}{\textcolor{blue}{$u$ many binary variables}}
\label{general:full-reformulation}
For a non-negative integer variable $x$ with upper bound $u$, modeled as 
\begin{equation}
0 \leq x \leq u, \ \ \ \ x \in \Z,
\end{equation}
this can be reformulated with $u$ binary variables $z_1, \dots, z_u$ as 
\begin{equation}
\begin{split}
x & = \sum_{i=1}^u i z_i = z_1 + 2 z_2 + \dots + u z_u\\
1 & \geq \sum_{i=1}^u z_i = z_1 + z_2 + \dots + z_u\\
z_i & \in \{0,1\} \ \ \text{ for } i=1, \dots, u
\end{split}
\end{equation}
\end{general}
We call this the \emph{full reformulation} because there is a binary variable $z_i$ associated with every value $i$ that $x$ could take.  That is, if $z_3 = 1$, then the second constraint forces $z_i = 0$ for all $i \neq 3$ (that is, $z_3$ is the only non-zero binary variable), and hence by the first constraint, $x = 3$.

\begin{general}{Binary reformulation}{\textcolor{blue}{$O(\log u)$ many binary variables}}
\label{general:log-reformulation}
For a non-negative integer variable $x$ with upper bound $u$, modeled as 
\begin{equation}
0 \leq x \leq u, \ \ \ \ x \in \Z,
\end{equation}
this can be reformulated with $u$ binary variables $z_1, \dots, z_{\lfloor\log( u )\rfloor}$ as 
\begin{equation}
\begin{split}
x & = \sum_{i=0}^{\lfloor\log( u )\rfloor}2^i z_i = z_0 + 2 z_1 +  4 z_2 + 8 z_3 + \dots + 2^{\lfloor\log( u )\rfloor} z_{\lfloor\log( u )\rfloor}\\
z_i & \in \{0,1\} \ \ \text{ for } i=1, \dots, \lfloor\log( u )\rfloor
\end{split}
\end{equation}
\end{general}
We call this the \emph{log reformulation} because this requires only logarithmically many binary variables in terms of the upper bound $u$.   This reformulation is particularly better than the full reformulation when the upper bound $u$ is a ``larger" number, although we will leave it ambiguous as to how larger a number need to be in order to be described as a ``larger" number. 

\section{Exercises}
\begin{ex}{Locating Distribution Terminals}{}{}

Suppose that a company based in St. John's is considering adding distribution terminals in (1) Halifax, (2) Moncton, (3) Montréal, (4) Ottawa, and (5) Toronto.\\

The cost of building the five terminals in millions of dollars would be 10 in Halifax, 12 in Moncton, 20 in Montréal, 18 in Ottawa, and 25 in Toronto. 
\begin{enumerate}
\item No more than two terminals may be built. 
\item If no terminal is built in Halifax, then one must be built in Moncton. 
\item At least one terminal must be built in central Canada (i.e. at locations (3), (4), or (5)).
\end{enumerate}

Define the variables:
$y_{i}= \begin{cases}1 & \text { if a terminal is built in city } i(i=1, \ldots, 5) \\ 0 & \text { otherwise }\end{cases}$

Using these variables, write a (binary) integer program that models the problem of determining which terimals to build subject to the constraints above that minimizes cost.
\end{ex}

\section{Literature and Resources}

\begin{resource}

\begin{itemize}
\item The AIMMS modeling has many great examples.  It can be book found here:\href{https://www.aimms.com/english/developers/resources/manuals/optimization-modeling}{AIMMS Modeling Book}.
\item \href{http://web.mit.edu/15.053/www/AMP-Chapter-09.pdf}{MIT Open Courseware}
\item  For many real world examples, see this book \href{https://link-springer-com.ezproxy.lib.vt.edu/book/10.1007/978-1-4939-1007-6}{
Case Studies in Operations Research
Applications of Optimal Decision Making, edited by  Murty, Katta G}.
Or find it \href{https://www.springer.com/gp/book/9781493910069}{here}.
\item \href{https://github.com/Gurobi/modeling-examples}{GUROBI modeling examples by GUROBI}
\item \href{https://github.com/open-optimization/open-optimization-or-examples/tree/master/integer-programming}{GUROBI modeling examples by Open Optimization that are linked in this book}
\end{itemize}

Knapsack Problem
\begin{itemize}
\item \href{https://www.youtube.com/watch?v=GteOoMGUOdY}{Video! - Michel Belaire (EPFL) teaching knapsack problem}
\end{itemize}

Set Cover
\begin{itemize}
\item \href{https://www.youtube.com/watch?v=cjSeHSjPmsk}{Video! - Michel Belaire (EPFL) explaining set covering problem}
\item See \href{https://download.aimms.com/aimms/download/manuals/AIMMS3OM_MediaSelection.pdf}{AIMMS - Media Selection} for an example of set covering applied to media selection.
\end{itemize}

Facility Location
\begin{itemize}
\item \href{https://en.wikipedia.org/wiki/Facility_location_problem}{Wikipedia - Facility Location Problem}
\item See \href{https://github.com/Gurobi/modeling-examples/tree/master/facility_location}{GUROBI Modeling Examples - Facility Location}.
\end{itemize}

{Other examples}
\begin{itemize}
\item \href{https://www.juliaopt.org/notebooks/JuMP-Sudoku.html}{Sudoku}
\item \href{https://download.aimms.com/aimms/download/manuals/AIMMS3OM_EmployeeTraining.pdf}{AIMMS - Employee Training}

\item \href{https://download.aimms.com/aimms/download/manuals/AIMMS3OM_MediaSelection.pdf}{AIMMS - Media Selection}

\item \href{https://download.aimms.com/aimms/download/manuals/AIMMS3OM_Diet.pdf}{AIMMS - Diet Problem}

\item \href{https://download.aimms.com/aimms/download/manuals/AIMMS3OM_FarmPlanning.pdf}{AIMMS - Farm Planning Problem}

\item \href{https://download.aimms.com/aimms/download/manuals/AIMMS3OM_Pooling.pdf}{AIMMS - Pooling Probem}

\item \href{https://www.informs.org/Impact}{INFORMS - Impact}
\item \href{https://www.informs.org/Impact/O.R.-Analytics-Success-Stories/Optimized-school-bus-routing-helps-school-districts-design-better-policies}{INFORMS - Success Story - Bus Routing}
\end{itemize}

{Notes from AIMMS modeling book.}
\begin{itemize}

\item \href{http://inside.mines.edu/~anewman/MIP_practice120212.pdf}{AIMMS - Practical guidelines for solving difficult MILPs}

\item \href{https://download.aimms.com/aimms/download/manuals/AIMMS3OM_LinearProgrammingTricks.pdf}{AIMMS - Linear Programming Tricks}

\item  \href{https://download.aimms.com/aimms/download/manuals/AIMMS3OM_FormulatingOptimizationModels.pdf}{AIMMS - Formulating Optimization Models}

\item \href{https://pdfs.semanticscholar.org/b01f/ad44c20c372fdda95cbfb980c0d37302de07.pdf}{AIMMS - Practical guidelines for solving difficult linear programs}
\end{itemize}

Modeling Tricks
\begin{itemize}
\item \href{https://jump.dev/JuMP.jl/stable/tutorials/linear/tips_and_tricks/}{JuMP tips and tricks}
\item \href{https://docs.mosek.com/modeling-cookbook/mio.html}{Mosek Modeling Cookbook}
\end{itemize}

{Further Topics}
\begin{itemize}
\item \href{https://or.stackexchange.com/questions/1319/best-model-for-precedence-constraints-within-scheduling-problem}{Precedence Constraints}
\end{itemize}
\end{resource}

\end{document}
